\documentclass[11pt,a4paper]{article}
\usepackage[margin=25mm]{geometry}
\usepackage{amsmath,amssymb,amsthm,hyperref}
\setlength{\emergencystretch}{2em}
\newtheorem{lemma}{Lemma}
\newcommand{\gen}[1]{\langle#1\rangle}
\title{Inheritance of perfect splits in actual wings}
\author{Human Compression project: a component of the L019 recovery}
\date{}
\begin{document}\maketitle
Let $G=\gen{\rho_0,\ldots,\rho_{r-1}}$ be transitive, with
string involutions and every maximal omission intransitive.
Suppose it has a perfect split at $i$, with left and right
orbits $A,B$, and write
\[
 q=\rho_i=(a\ b)\alpha\beta,
 \qquad a\in A,quad b\in B.
\]
The marks before $i$ are supported and transitive on $A$,
and those after it are supported and transitive on $B$.
The internal factors $\alpha,\beta$ fix $a,b$ respectively.
A one-point wing has its empty transitive marking.

\begin{lemma}\label{omissions}
Every maximal omission of either actual wing is intransitive
on that wing's support.
\end{lemma}
\begin{proof}
If a wing remained transitive after deletion of one of its
marks, the other full wing and the crossing mark $q$ would
join its support to the other support. The corresponding
maximal omission of $G$ would then be transitive, a contradiction.
\end{proof}

\begin{lemma}[A nonadjacent wing split]\label{far}
Suppose the right wing has a perfect split at $j\ge i+2$.
Then $j$ is a perfect split of the parent marking, even if
$j=r-1$. Either trivial internal factor remains trivial.
The reversed statement holds for a left-wing split at $j\le i-2$,
including $j=0$.
\end{lemma}
\begin{proof}
Write the wing's two sides as $C,D$, where the marks strictly
between $i$ and $j$ are transitive and supported on $C$,
and those after $j$ are transitive and supported on $D$.
Their unique crossing edge is $\{c,d\}$, with $c\in C$,
$d\in D$. Every mark after $i+1$ fixes $b$ by the original
unique crossing. The mark $\rho_{i+1}$ moves $b$, since
the right wing is transitive. As it lies before $j$, its
support is contained in $C$, so $b\in C$.

The parent omission $G_j$ contains groups transitive on
$A,C,D$ separately, and $q$ joins $A$ to $C$ through
$\{a,b\}$. If $q$ had any transposition from $C$ to $D$,
this omission would be transitive on the full set. Thus $q$
preserves $D$ and its complement setwise. It commutes with
$\rho_j$, and $\{c,d\}$ is the only $\rho_j$-edge between
these two sets. Hence $q$ preserves this edge and fixes each
of its endpoints. It also centralizes the transitive group
on $D$ generated by the marks after $j$, so fixes $D$
pointwise. When $D$ is a singleton, preservation already
gives the same conclusion; no nontrivial-tail hypothesis
is needed.

The marks before $j$ therefore generate a transitive group
on $A\cup C$ and fix $D$, while the marks after $j$ remain
transitive and supported on $D$. The two actual omission
orbits are $A\cup C,D$. The $j$-mark fixes $A$, so its
internal factors are unchanged apart from adjoining fixed
points to the first side. This proves both conclusions.
\end{proof}

\begin{lemma}[The adjacent split of a sufficiently large wing]\label{near}
Let the right wing have degree $m=|B|$ and rank $t=r-i-1$.
Suppose $m\ge2$, $2t>m-1$, and its first index $j=i+1$
is a perfect split. Then $j$ is a perfect split of the parent
marking and its internal factor on the new left side is trivial.
In addition, the original factor $\beta$ is trivial.
\end{lemma}
\begin{proof}
The first wing split has a singleton near side and a far
side $D$ on which $T=G_{>j}$ is transitive. Its singleton
is $b$. Indeed $T$ fixes $b$, while only the first wing
mark can move it. If $|D|>1$, the singleton is the only
$T$-fixed point in $B$; if $|D|=1$, choose the side names
so that the moved point $b$ is the near side. Thus
$B=\{b\}\sqcup D$ and $\rho_j=(b\ d)\delta$, where
$\delta$ is supported on $D\setminus\{d\}$.

The factor $\beta$ is supported on $D$, centralizes $T$,
and, if nonidentity, is fixed-point-free there. Its pairs
would form a $T$-invariant system of $|D|/2$ blocks.
Choose a subfamily $M$ of the specified $T$ marks transitive
on this pair set, with at most $|D|/2-1$ members, using the
acquired minimal-transitivity forest. On $D$, the group
$\gen{M,q}$ contains the pair flips supplied by $q|_D=\beta$
and is transitive: it is transitive on the pairs and each
pair can be interchanged. Add all left-wing marks and
$\rho_j$. The resulting group is transitive on $A$, joins
$A$ to $b$ through $q$, and joins $b$ to $D$ through
$\rho_j$, hence is transitive on the whole domain.
Any $T$ mark absent from $M$ would contradict the parent
omission hypothesis. Therefore
\[
 t-1\le |M|\le |D|/2-1,
 \qquad\hbox{so}\qquad 2t\le |D|=m-1,
\]
contrary to the hypothesis. Thus $\beta=1$.

The marks through $i$ now act transitively on $A\cup\{b\}$
and fix $D$, while $T$ is transitive on $D$ and fixes the
new left side. The mark $\rho_j$ has only its crossing
$\{b,d\}$ on the new left side and fixes all of $A$.
It is therefore a perfect split with trivial left internal
factor. This includes a singleton $D$, where $\beta=1$
was automatic. Reversal gives the corresponding left-wing
assertion.
\end{proof}

\paragraph{Attribution and remaining scope.}
These are the actual support mechanisms behind the inheritance
steps in Cameron--Fernandes--Leemans, \emph{The number of string
C-groups of high rank}, arXiv:2212.12723v2, Section~8.
The nonadjacent proof makes explicit the parent omission
hypothesis and the terminal singleton case. The adjacent
proof replaces a classification bound by minimal transitivity
on actual involution pairs; it assumes no quotient IP.
It also strengthens the usable interface to the stated wing
rank condition. Neither lemma claims that an available wing
split is internal, or that the full stability correspondence
has been recovered.
\end{document}
